\section{Estimator internals and the meaning of a portable event study}
\label{sec:exec28-panel-estimator-contracts}

A branch-policy comparison can survive a file conversion while losing its
scientific meaning. A bank that moves a servicing-policy panel from Stata into
an R wrapper must preserve treatment codes, missingness reasons, observation
clocks, the population contributing to each contrast and the origin of every
interval. A common column layout is useful only after these distinctions are
made explicit. The following comparison uses nine additional files at the
previously fixed \texttt{fastdid} and \texttt{nonabsdid} commits. It is a static
implementation review, accompanied by independent arithmetic, and does not
claim that either package or its imported numerical dependencies was executed
\citep{exec28_panel_estimator_sources,exec28_panel_estimator_reference}.

\paragraph{Effective support and nuisance estimation.}
In the selected balanced \texttt{fastdid} estimator, the control contrast uses
normalized propensity odds when covariates and the chosen IPW or doubly robust
option activate that branch. Outcome adjustment fits a weighted linear model
to the control units' outcome changes. The code includes nuisance-estimation
terms in its influence calculation; thus an adapter that exports only a
coefficient and discards the fitted weighting law cannot reconstruct the same
uncertainty calculation. The unbalanced branch instead forms separately
normalized pre- and post-period means, weights propensity fitting by presence
in those periods and fits separate outcome regressions. Those branches need
distinct declared observation laws and tests. Their shared function name is
not a proof of equivalent targets or calibration.

The selected balanced branch warns about extreme fitted propensities and caps
them at $1-10^{-10}$. The corresponding control odds can still be
$9{,}999{,}999{,}999$ before sampling-weight multiplication. A warning and a
finite numerical value therefore provide no useful-overlap guarantee. That
branch also replaces missing or magnitude-above-$10^{10}$ entries of its
scaled covariance matrix by zero. The proposal must expose this alteration,
retain the original diagnostic and require a declared refusal or regularization
policy; it cannot quietly interpret zeroed entries as established certainty.
The selected unbalanced branch does not contain the same replacement at its
corresponding covariance step. Dependency versions, convergence and singular
design behavior remain separate qualification questions.

For weights $w_i$, the source calculates the effective support
\[
 n_{\mathrm{eff}}=\frac{(\sum_i w_i)^2}{\sum_i w_i^2}.
\]
It rejects a group, or an unbalanced group-period, when this quantity is
nonfinite or below two. Two equally weighted observations give
$n_{\mathrm{eff}}=2$; weights $(1,2)$ give $9/5$ and trigger the selected
refusal despite there being two physical observations. The residual component
is multiplied by $\sqrt{n_{\mathrm{eff}}/(n_{\mathrm{eff}}-1)}$ after this
check. For two equal weights its squared multiplier is two. This algebra
describes the implemented adjustment, not a general finite-sample coverage
theorem for estimated weights, correlated branches or selected cells. A public
comparison should report physical counts, effective counts, refused cells and
surviving aggregation weights separately. The upstream propensity solver,
weighted least-squares behavior and multiplier bootstrap still require their
own pinned review and matched execution contract.

\paragraph{Calendar and argument conversion.}
The auxiliary \texttt{fastdid} path normalizes uniformly spaced times and
rejects irregular spacing. Calendar-valued group and time outputs are mapped
back, while dynamic event-time coordinates remain in normalized period units.
Sampling weights and time-invariant inputs are taken from the first unit row
in the unbalanced auxiliary construction. A bank must therefore define which
period's exposure weight represents a unit, and must not silently treat a
one-period event coordinate as one calendar day. Custom aggregation expressions
also require an explicitly controlled interface. The proposed CLI/MCP should
expose named, validated aggregation choices rather than pass arbitrary business
text into the package's expression option.

The \texttt{nonabsdid} alias layer rejects simultaneous alias and canonical
arguments. It maps \texttt{effects} to \texttt{leads=effects-1} and
\texttt{placebo} to \texttt{lags}; its selected validation checks numeric range
but then uses integer coercion. At that layer, an input of $2.9$ effects becomes
two effects and one lead, while $1.9$ placebos becomes one lag. These are
source-level conversion witnesses, not native invocation results. The proposed
boundary should require finite integer horizons before dispatch and return
both requested and resolved options. The DCDH tidier subtracts one from the
upstream time coordinate and reads estimates from embedded plot data. A call
that omits the plot is refused at the inspected extraction boundary. Preserve
the native clock alongside the display clock and qualify the exact upstream
object version instead of treating a visible plot slot as a permanent API.

\paragraph{Intervals, reference rows and file round trips.}
The \texttt{fect} tidier carries native interval bounds and takes standard
errors only when the designated matrix has a matching number of rows. This
length check does not independently establish that its row identities match
the event-time vector. Its method-label helper maps missing and unrecognized
labels to \texttt{IFE}; an integration should retain the original label and
flag an unknown method rather than infer that a named estimator ran. The
PanelMatch tidier requires named estimate coordinates, handles separate
post-treatment and placebo fields and can append a zero-estimate, zero-error
row at time $-1$. That row is a plotting reference, not an observed exact-zero
effect, and must be excluded from empirical coverage and precision denominators.
Existing native bounds, normal approximations and simultaneous bands must keep
distinct provenance through the common constructor; that constructor and
upstream versions remain explicit follow-up inspections.

The selected Stata reader defaults to converting labelled values into factors
and collapsing tagged missing values to ordinary missingness. Those defaults
may be convenient for presentation but are not a reversible preservation of a
numeric treatment code or a reason-for-missingness taxonomy. Preprocessing must
retain those mappings before fitting. The writer sanitizes names, truncates
them to 32 characters and then makes them unique. Two names sharing the same
32-character prefix consequently produce a 34-character second name after
the selected suffix operation. This is a pre-writer boundary witness; whether
the exact imported writer rejects or transforms it has not been executed.
A portable schema should allocate collision-safe names within the length bound
and preserve a reverse mapping.

These findings support a conditional adapter design: preserve raw and resolved
arguments, native and displayed clocks, treatment labels, missingness tags,
effective samples, influence provenance and reference-row status. They do not
support an unconditional port or adoption decision. Full upstream numerical
and constructor inspection, exact dependency and rights profiles, the broader
multiple-event paper edition, and native public fixtures remain prerequisites
to stronger claims. Unknown counterfactual trends remain scientific assumptions
that a faithful file or API conversion cannot establish.
